\begin{pastexam}[label={px:ust-115-ee-ds-02}]{臺聯大電機類 115 學年度 第 2 題}
（單選；2 分）

What are the main disadvantages of storing a sparse matrix using a general two-dimensional array?
\begin{enumerate}[label=(\Alph*)]
\item Access speed is too slow
\item The code is too complex
\item Wasting a lot of memory space
\item Matrix multiplication cannot be performed
\end{enumerate}
\end{pastexam}

\begin{pastsoln}
答案是 (C)。

稀疏矩陣是大部分元素為 $0$ 的矩陣。 $m \times n$ 的矩陣只有 $t$ 個非零元素（$t \ll mn$）時，
用一般二維陣列存仍然要 $\Theta(mn)$ 的空間，絕大多數格子都在存 $0$。 改用三元組
（列、行、值）只存非零元素，空間是 $\Theta(t)$（資結書第~4.3~節）。 例如 $1000 \times 1000$
的矩陣只有 $2000$ 個非零元素，二維陣列要存一百萬個數，三元組只要 $3 \times 2000 = 6000$ 個。
除了空間，轉置、相乘時也會花很多時間在處理 $0$ 上。

逐項檢查：
\begin{itemize}
\item (A) 二維陣列用位址公式直接算出任何一格的位置，存取是 $\Theta(1)$，反而是最快的。 錯。
\item (B) 二維陣列的程式最簡單，三元組的轉置、相乘才複雜。 錯。
\item (C) 對，見上。
\item (D) 二維陣列當然可以做矩陣乘法，只是會做很多乘以 $0$ 的白工。 錯。
\end{itemize}
\end{pastsoln}
